\subsection{Procedures as General Methods}
\label{Section 1.3.3}

We introduced compound procedures in \link{Section 1.1.4} as a mechanism for
abstracting patterns of numerical operations so as to make them independent of
the particular numbers involved.  With higher-order procedures, such as the
\code{integral} procedure of \link{Section 1.3.1}, we began to see a more
powerful kind of abstraction: procedures used to express general methods of
computation, independent of the particular functions involved.  In this section
we discuss two more elaborate examples---general methods for finding zeros and
fixed points of functions---and show how these methods can be expressed
directly as procedures.

\subsubsection*{Finding roots of equations by the half-interval method}

The \newterm{half-interval method} is a simple but powerful technique for
finding roots of an equation \( f(x) = 0 \), where \( f \) is a continuous
function.  The idea is that, if we are given points \( a \) and \( b \) such that
\( f(a) < 0 < f(b) \), then \( f \) must have at least one zero between
\( a \) and \( b \).  To locate a zero, let \( x \) be the average of \( a \) and \( b \),
and compute \( f(x) \).  If \( f(x) > 0 \), then \( f \) must have a zero
between \( a \) and \( x \).  If \( f(x) < 0 \), then \( f \) must have a zero
between \( x \) and \( b \).  Continuing in this way, we can identify smaller and
smaller intervals on which \( f \) must have a zero.  When we reach a point where
the interval is small enough, the process stops.  Since the interval of
uncertainty is reduced by half at each step of the process, the number of steps
required grows as \( \Theta(\log(L / T)) \), where \( L \) is the
length of the original interval and \( T \) is the error tolerance (that is, the
size of the interval we will consider ``small enough'').  Here is a procedure
that implements this strategy:

\begin{codeBlock}
>>> def search(f, neg_point, pos_point):
...     midpoint = average(neg_point, pos_point)
...     if is_close_enough(neg_point, pos_point):
...         return midpoint
...     test_value = f(midpoint)
...     if test_value > 0:
...         return search(f, neg_point, midpoint)
...     elif test_value < 0:
...         return search(f, midpoint, pos_point)
...     else:
...         return midpoint
\end{codeBlock}

\noindent
We assume that we are initially given the function \( f \) together with points
at which its values are negative and positive.  We first compute the midpoint
of the two given points.  Next we check to see if the given interval is small
enough, and if so we simply return the midpoint as our answer.  Otherwise, we
compute as a test value the value of \( f \) at the midpoint.  If the test value
is positive, then we continue the process with a new interval running from the
original negative point to the midpoint.  If the test value is negative, we
continue with the interval from the midpoint to the positive point.  Finally,
there is the possibility that the test value is 0, in which case the midpoint
is itself the root we are searching for.

To test whether the endpoints are ``close enough'' we can use a procedure
similar to the one used in \link{Section 1.1.7} for computing square
roots:\footnote{We have used 0.001 as a representative ``small'' number to
indicate a tolerance for the acceptable error in a calculation.  The
appropriate tolerance for a real calculation depends upon the problem to be
solved and the limitations of the computer and the algorithm.  This is often a
very subtle consideration, requiring help from a numerical analyst or some
other kind of magician.}

\begin{codeBlock}
>>> def is_close_enough(x, y):
...     return abs(x - y) < 0.001
\end{codeBlock}

\noindent
\code{search} is awkward to use directly, because we can accidentally give it
points at which \( f \)'s values do not have the required sign, in which case we
get a wrong answer.  Instead we will use \code{search} via the following
procedure, which checks to see which of the endpoints has a negative function
value and which has a positive value, and calls the \code{search} procedure
accordingly.  If the function has the same sign on the two given points, the
half-interval method cannot be used, in which case the procedure signals an
error.\footnote{The original does this with \code{error}, which prints its
arguments as a message.  Python's equivalent is to \newterm{raise} an
exception, and the argument to \code{ValueError} is the message.  We shall
have more to say about exceptions when we need them; for now it is enough that
raising one stops the computation and reports why.}

\begin{codeBlock}
>>> def half_interval_method(f, a, b):
...     a_value, b_value = f(a), f(b)
...     if a_value < 0 and b_value > 0:
...         return search(f, a, b)
...     elif b_value < 0 and a_value > 0:
...         return search(f, b, a)
...     else:
...         raise ValueError(
...             f"Values are not of opposite sign: {a}, {b}")
\end{codeBlock}

\noindent
The following example uses the half-interval method to approximate \( \pi \) as
the root between 2 and 4 of \( \sin x = 0 \):

\begin{codeBlock}
>>> import math
>>> half_interval_method(math.sin, 2.0, 4.0)
~\textit{3.14111328125}~
\end{codeBlock}

\noindent
Here is another example, using the half-interval method to search for a root of
the equation \( x^3 - 2x - 3 = 0 \) between 1 and 2:

\begin{codeBlock}
>>> half_interval_method(lambda x: x * x * x - 2 * x - 3,
...                       1.0,
...                       2.0)
~\textit{1.89306640625}~
\end{codeBlock}

\subsubsection*{Finding fixed points of functions}

A number \( x \) is called a \newterm{fixed point} of a function \( f \) if \( x \)
satisfies the equation \( f(x) = x \).  For some functions \( f \) we can
locate a fixed point by beginning with an initial guess and applying \( f \)
repeatedly,

$$ f(x),\quad f(f(x)),\quad f(f(f(x))), \quad\dots, $$

\noindent
until the value does not change very much.  Using this idea, we can devise a
procedure \code{fixed\_point} that takes as inputs a function and an initial
guess and produces an approximation to a fixed point of the function.  We apply
the function repeatedly until we find two successive values whose difference is
less than some prescribed tolerance:

\begin{codeBlock}
>>> tolerance = 0.00001

>>> def fixed_point(f, first_guess):
...     def is_close_enough(v1, v2):
...         return abs(v1 - v2) < tolerance
...     def try_guess(guess):
...         next_guess = f(guess)
...         if is_close_enough(guess, next_guess):
...             return next_guess
...         else:
...             return try_guess(next_guess)
...     return try_guess(first_guess)
\end{codeBlock}

\noindent
For example, we can use this method to approximate the fixed point of the
cosine function, starting with 1 as an initial approximation:\footnote{Try this
during a boring lecture: Set your calculator to radians mode and then
repeatedly press the \code{cos} button until you obtain the fixed point.}

\begin{codeBlock}
>>> fixed_point(math.cos, 1.0)
~\textit{0.7390822985224023}~
\end{codeBlock}

\noindent
Similarly, we can find a solution to the equation
\( y = \sin y + \cos y \):

\begin{codeBlock}
>>> fixed_point(lambda y: math.sin(y) + math.cos(y),
...             1.0)
~\textit{1.2587315962971173}~
\end{codeBlock}

\noindent
The fixed-point process is reminiscent of the process we used for finding
square roots in \link{Section 1.1.7}.  Both are based on the idea of repeatedly
improving a guess until the result satisfies some criterion.  In fact, we can
readily formulate the square-root computation as a fixed-point search.
Computing the square root of some number \( x \) requires finding a \( y \) such
that \( y^2 = x \).  Putting this equation into the equivalent form
\( y = x / y \), we recognize that we are looking for a fixed point of the
function\footnote{\( \mapsto \) (pronounced ``maps to'') is the mathematician's way of
writing \code{lambda}.  \( y \mapsto x / y \) means \code{lambda y: x / y},
that is, the function whose value at \( y \) is \( x / y \).} \( y \mapsto x / y \),
and we can therefore try to compute square roots as

\begin{codeBlock}
>>> def sqrt(x):
...     return fixed_point(lambda y: x / y, 1.0)
\end{codeBlock}

\noindent
Unfortunately, this fixed-point search does not converge.  Consider an initial
guess \( y_1 \).  The next guess is \( y_2 = x / y_1 \) and the next guess is
\( y_3 = x / y_2 = x / (x / y_1) = y_1 \).  This results in an
infinite loop in which the two guesses \( y_1 \) and \( y_2 \) repeat over and
over, oscillating about the answer.

One way to control such oscillations is to prevent the guesses from changing so
much.  Since the answer is always between our guess \( y \) and \( x / y \), we
can make a new guess that is not as far from \( y \) as \( x / y \) by averaging
\( y \) with \( x / y \), so that the next guess after \( y \) is
\( \frac{1}{2}(y + x / y) \)
instead of \( x / y \).  The process of making such a sequence of
guesses is simply the process of looking for a fixed point of
\( y \mapsto \frac{1}{2}(y + x / y) \):

\begin{codeBlock}
>>> def sqrt(x):
...     return fixed_point(lambda y: average(y, x / y), 1.0)
\end{codeBlock}

\noindent
(Note that \( y = \frac{1}{2}(y + x / y) \) is a simple transformation of the
equation \( y = x / y; \) to derive it, add \( y \) to both sides of the
equation and divide by 2.)

With this modification, the square-root procedure works.  In fact, if we
unravel the definitions, we can see that the sequence of approximations to the
square root generated here is precisely the same as the one generated by our
original square-root procedure of \link{Section 1.1.7}.  This approach of
averaging successive approximations to a solution, a technique that we call
\newterm{average damping}, often aids the convergence of fixed-point searches.

\begin{exerciseGroup}
\phantomsection\label{Exercise 1.35}
\exerciseHeading{Exercise 1.35:} Show that the golden ratio
\( \varphi \) (\link{Section 1.2.2}) is a fixed point of the transformation
\( x \mapsto 1 + 1 / x \), and use this fact to compute \( \varphi \) by means
of the \code{fixed\_point} procedure.

\phantomsection\label{Exercise 1.36}
\exerciseHeading{Exercise 1.36:} Modify \code{fixed\_point} so that
it prints the sequence of approximations it generates, using the \code{newline}
and \code{display} primitives shown in \link{Exercise 1.22}.  Then find a
solution to \( x^x = 1000 \) by finding a fixed point of \( x \mapsto
\log(1000) / \log(x) \).  (Use Scheme's primitive \code{log}
procedure, which computes natural logarithms.)  Compare the number of steps
this takes with and without average damping.  (Note that you cannot start
\code{fixed\_point} with a guess of 1, as this would cause division by
\( \log(1) = 0 \).)

\phantomsection\label{Exercise 1.37}
\exerciseHeading{Exercise 1.37:} \begin{enumerate}[a.]

\item
An infinite \newterm{continued fraction} is an expression of the form

$$ {f} = \cfrac{N_1}{D_1 + \cfrac{N_2}{D_2 + \cfrac{N_3}{D_3 + \dots}}}\,. $$

As an example, one can show that the infinite continued fraction expansion with
the \( N_i \) and the \( D_i \) all equal to 1 produces \( 1 / \varphi \), where
\( \varphi \) is the golden ratio (described in \link{Section 1.2.2}).  One way to
approximate an infinite continued fraction is to truncate the expansion after a
given number of terms.  Such a truncation---a so-called \newterm{\textit{k}-term
finite continued fraction}---has the form

$$ \cfrac{N_1}{D_1 + \cfrac{N_2}{\ddots + \cfrac{N_k}{D_k}}}\,. $$

Suppose that \code{n} and \code{d} are procedures of one argument (the term
index \( i \)) that return the \( N_i \) and \( D_i \) of the terms of the
continued fraction.  Define a procedure \code{cont\_frac} such that evaluating
\code{cont\_frac(n, d, k)} computes the value of the \( k \)-term finite
continued fraction.  Check your procedure by approximating \( 1 / \varphi \) using

\begin{codeBlock}
cont_frac(lambda i: 1.0,
          lambda i: 1.0,
          k)
\end{codeBlock}

\noindent
for successive values of \code{k}.  How large must you make \code{k} in order
to get an approximation that is accurate to 4 decimal places?

\item
If your \code{cont\_frac} procedure generates a recursive process, write one
that generates an iterative process.  If it generates an iterative process,
write one that generates a recursive process.

\end{enumerate}

\phantomsection\label{Exercise 1.38}
\exerciseHeading{Exercise 1.38:} In 1737, the Swiss mathematician
Leonhard Euler published a memoir \textit{De Fractionibus Continuis}, which
included a continued fraction expansion for \( e - 2 \), where \( e \) is the base
of the natural logarithms.  In this fraction, the \( N_i \) are all 1, and
the \( D_i \) are successively 1, 2, 1, 1, 4, 1, 1, 6, 1, 1, 8, \( \dots \).
Write a program that uses your \code{cont\_frac} procedure from \link{Exercise 1.37}
to approximate \( e \), based on Euler's expansion.

\phantomsection\label{Exercise 1.39}
\exerciseHeading{Exercise 1.39:} A continued fraction
representation of the tangent function was published in 1770 by the German
mathematician J.H. Lambert:

$$ {\tan x} = \cfrac{x}{1 - \cfrac{x^2}{3 - \cfrac{x^2}{5 - \dots}}}\,, $$

\noindent
where \( x \) is in radians.  Define a procedure \code{tan\_cf(x, k)} that
computes an approximation to the tangent function based on Lambert's formula.
\code{k} specifies the number of terms to compute, as in \link{Exercise 1.37}.
\end{exerciseGroup}

